% Modernised reading — fonds Alexandre Grothendieck, shelfmark 118, pages 1-40
% (the whole folder).
%
% This is an INTERPRETATION, not a transcription. It re-reads the pages in
% today's notation and vocabulary, states the hypotheses the manuscript leaves
% implicit, and drops the critical apparatus so that the mathematics reads
% straight through. Everything the brackets and underlines carried moves into
% a footnote. It is held to being mathematically correct as it stands; where
% the manuscript is loose or mistaken, the true statement is what appears here
% and the footnote says what the page has. In any dispute of substance the
% transcription governs: whoever cites Grothendieck must cite
% batch-01.fr.tex and batch-02.fr.tex.
%
% Pass: Opus 5.5 (claude-opus-5-5), 2026-10-03 — first pass, unchecked against the pages by a human.
%
% Scope: the folder was read whole — two batches, pages 1-40 — before any of
% this was written, and it is ONE file for the shelfmark. Pages 3, 5, 6, 9,
% 12 and 17 carry nothing of his mathematics (accounts, administrative
% tables, an envelope, a sheet marked « Poèmes », a typed notice dated
% 29 April 1985) and page 37 is blank; they are not transcribed and are
% mentioned only in the spine section.
%
% Notation fixed for the whole document.
%  - \mathcal W is the localisateur fondamental of (Cat) — his underlined W
%    on pp. 14-40, his plain W on pp. 2-8 and in places after; read as one
%    class throughout. \mathcal W_A its trace on the presheaves \widehat A,
%    \mathcal W_0 the usual weak equivalences (p. 35).
%  - Class letters kept as his, in roman: F, C (Grothendieck fibrations and
%    cofibrations = opfibrations, NOT model-category ones), Sm, Pr, S, B, T,
%    WSF, WSC, WSB, As, As°, Tot As.
%  - X_{/b} and {}_{b\backslash}X the comma categories of f over and under b,
%    X_b the fibre; « ouvert » = crible (sieve), « fermé » = cocrible.
%  - \int(f_0, f_1) the Grothendieck construction of a span (his ∫ of a
%    matrix); C(f) the cylinder of p. 10, C(Y) the cone with final vertex,
%    C'(Y) with initial vertex.
%  - Hot = Hot(\mathcal W) = \mathcal W^{-1}(Cat), Hot_A, \gamma the
%    localisation functor; \overline{Cat}, \overline{\widehat A} the
%    categories of homotopy classes (his barred objects).
%
% Substantive departures from the pages (footnoted where they occur):
%  p. 2/4 As = W ∩ WSC and S ∩ W = S ∩ T kept as containments where only the
%  containment is evident (the converse is a Theorem-B statement); p. 8 the
%  box « WSC ∩ Pr = S » read ⊂ S; p. 8 class letter read C, not E; p. 13
%  f^♯ identified with the open immersion f^♮; p. 14 the « cart » square of
%  (4) read as cocartesian; p. 14 the first Y' of « 4 Thm » read Y; p. 16
%  h^{-1}(Y') read h^{-1}(Y); p. 20 the square supposed cartesian (the
%  fibre formula Y_{x'} = Y'_{x'} × X_{x'} needs it), and the Corollary's
%  derivation supplied; p. 22 « π_1 ≠ 1 » is not enough for the example —
%  take a group with non-zero abelianisation (an acyclic group gives a
%  contractible suspension), and the example disproves only g ∈ W ⇒ f ∈ W;
%  p. 22 the two cones read with initial vertex; p. 26 « fractions à
%  droite » (against p. 23's « à gauche ») and « lim finies » read as the
%  cospan calculus and finite colimits; p. 28 « copies de X_n » read as
%  copies of the elementary interval; p. 30 the relation s_j i_j = id
%  footnoted; pp. 32-34 the Lemma's « conditions équivalentes » is false as
%  stated and is restated as what pp. 33-35 prove; p. 38 L 3 needed in a
%  form with a natural transformation (f(X_i) ⊂ Y_i gives only ψf ≤ φ), and
%  the reduction by duality needs W = W°.
%
% Announced and not established: duality W = W° (p. 14, item 1); the
% questions of p. 4 (*) and p. 8 (« ⟸ ? »); whether the base-change square
% of p. 19 is cartesian (« À vérifier ! »); the unfinished Lemme of p. 20;
% the 1-equivalence surmise of p. 22; the class 𝔐 and the calculus of
% fractions in (Cat) (p. 23, « itou dans Cat »), which the proof of p. 40
% uses; Corollary 1 for infinite sums (p. 26) — surjectivity argued on
% pp. 29-30, injectivity reduced to a Kan-type condition and left open; the
% Lemma of p. 30, sketched only; infinite products (pp. 33, 35-36); the
% correspondence functors X → I / filtrations (« Il faudra expliciter »,
% p. 39); the « Cor ? » of p. 39; the Cor « W_A est saturé » (p. 39),
% stated without proof (one line, supplied here).
\documentclass[11pt,a4paper]{article}
\input{../preamble/grothendieck}

\folder{118}
\pages{1}{40}
\dating{[à partir de 1983]}
\watermark{Édition de démonstration\\interprétation personnelle de l'œuvre}
\foldertitle{Cohomology properties of maps in (Cat) : notes manuscrites (s.d.) — lecture modernisée du dossier entier}

\begin{document}

\section*{Résumé}

\begin{resume}
Une petite catégorie peut se regarder comme un espace : on la remplace par son
nerf, un objet combinatoire qui se réalise en espace topologique, et deux
catégories sont « faiblement équivalentes » quand ces espaces ont le même type
d'homotopie. Ces quarante pages, en anglais puis en français, prennent ce point
de vue au sérieux et cherchent, parmi les foncteurs entre petites catégories,
les analogues des applications qu'un topologue sait manier : fibrations,
applications propres et lisses, revêtements, immersions ouvertes et fermées,
cylindres et cônes.

Le problème vient de ce que les opérations naturelles sur les catégories
respectent mal l'homotopie. Recoller deux catégories le long d'une troisième
--- une somme amalgamée --- ne donne pas en général la catégorie dont le nerf
serait le recollement des nerfs ; inverser les équivalences faibles produit une
catégorie « des types d'homotopie » dont on ne sait pas d'emblée si elle a des
sommes ou des produits, ni si une flèche qui y devient inversible était déjà
une équivalence faible. Le dossier s'attaque à ces trois questions, et il le
fait pour une notion d'équivalence faible donnée par axiomes, ce que
Grothendieck nomme un localisateur fondamental : une classe $\mathcal{W}$ de
foncteurs stable par composition, contenant les catégories à objet final
parmi les « asphériques », et qui se teste fibre à fibre.

L'idée directrice est empruntée à la géométrie algébrique. Là, un morphisme
propre ou lisse est celui pour lequel la cohomologie se comporte bien par
changement de base, et une fibration de Serre est celle dont les fibres
déterminent la topologie. Ici, une catégorie au-dessus d'une autre se lit par
ses catégories « au-dessus » et « au-dessous » de chaque objet de la base ; on
classe les foncteurs selon la façon dont ces catégories varient, et le
classement se dispose en treillis, que les premières pages dessinent. Puis
vient l'outil qui rend les recollements maniables : une immersion ouverte ou
fermée, c'est-à-dire l'inclusion d'un morceau qu'aucune flèche ne quitte, ou
dans lequel aucune n'entre, permet de décrire une catégorie par ses deux
morceaux et la manière dont les flèches vont de l'un à l'autre. Avec lui, une
somme amalgamée le long d'une immersion transporte les équivalences faibles,
sous des hypothèses de fibration --- et le dossier montre par un exemple qu'on
ne peut pas les lever.

Le reste en tire les conséquences. Dans une catégorie de préfaisceaux, toute
inclusion se prête au recollement, et cela donne un calcul de fractions : toute
flèche de la catégorie homotopique s'écrit comme une flèche suivie de
l'inverse d'une équivalence faible. On en déduit les sommes finies, puis les
produits finis sous une hypothèse sur la catégorie de base que le dossier
dégage et analyse, une stabilité par réunions filtrées, et enfin que
$\mathcal{W}$ est saturé : une flèche est une équivalence faible dès qu'elle
devient inversible. Plusieurs questions restent ouvertes, que les feuillets
posent avec franchise : les sommes infinies (« C'est gênant = bien
gênant ! »), les produits infinis, l'invariance de $\mathcal{W}$ par passage
à la catégorie opposée.

Les noms modernes à chercher sont ceux de la théorie des localisateurs
fondamentaux, développée ensuite par Cisinski et Maltsiniotis, de la structure
de modèles de Thomason sur les petites catégories, des théorèmes A et B de
Quillen, et des catégories test de \emph{Pursuing Stacks}, qu'il écrit à la
même époque.

\keywords{basic localizer, Grothendieck fibration, opfibration, smooth functor, proper functor, Quillen's Theorem A, Quillen's Theorem B, sharp map, comma category, mapping cylinder, sieve, cosieve, collage of a profunctor, Grothendieck construction, homotopy pushout, Thomason model structure, Dwyer map, local test category, totally aspherical category, calculus of fractions, category of cofibrant objects, strong saturation}
\end{resume}

\section*{Le fil du dossier, et les conventions}

Montpellier date le dossier « à partir de 1983 » et le range parmi les papiers
qui entourent \emph{Pursuing Stacks}. Seul un feuillet est daté, la page 17,
un avis de soutenance de thèse du 29 avril 1985 sans rapport avec les notes ;
les autres feuillets sans mathématiques (pages 3, 5, 6, 9 et 12 : des comptes,
des tableaux administratifs, une enveloppe, une feuille portant « 1976 » et
« Poèmes ») et la page blanche 37 sont laissés de côté. Une chemise rose
(page 1) porte le titre, « Cohomology properties of maps in (Cat) », et la
première rubrique, « 1) Fibration-like properties »\footnote{« like » est lu
sans certitude. La page 10 ouvre une rubrique « 2) », dont le titre, encadré
et en partie rayé, se lit avec doute « Immersions fermées et ouvertes,
propriétés cohomologiques ».}. Les pages 18 à 25 forment une suite numérotée
par lui de 1 à 7 (« 1 », « 2 », « 2 bis », puis 3 à 7) ; au-delà, aucun
numéro de sa main ne se voit.

\begin{enumerate}
\item[1.] \emph{Classes de foncteurs} (pages 2 à 8) : treillis de classes ---
fibrations, cofibrations, lisses, propres, Serre, leurs variantes faibles et
leurs intersections avec $\mathcal{W}$ --- et leurs caractérisations par les
catégories comma.
\item[2.] \emph{Factorisations, immersions, cylindres} (pages 10 à 13).
\item[3.] \emph{Le programme} (pages 14 et 15), qui annonce les stations
suivantes, numérotées de (0) à (6).
\item[4.] \emph{Recollements} (page 16) : une catégorie décrite par un ouvert,
son fermé complémentaire et un profoncteur.
\item[5.] \emph{Sommes amalgamées le long d'une immersion} (pages 18 à 22) :
le théorème dans (Cat), sa version dans les préfaisceaux, et le
contre-exemple.
\item[6.] \emph{La catégorie homotopique comme catégorie de fractions} (pages
23 à 25).
\item[7.] \emph{Sommes} (pages 26 à 30) et \emph{produits finis} (pages 31 à
36) dans $\mathrm{Hot}$.
\item[8.] \emph{Filtrations} (pages 38 et 39) et \emph{saturation} (pages 39
et 40).
\end{enumerate}

Le programme de la page 14 est réalisé presque point par point : ses rubriques
(3) et (4) aux pages 18 à 22, sa rubrique (5) aux pages 23 à 36, sa rubrique
(6) aux pages 39 et 40. Sa rubrique (1), la dualité, n'est pas traitée, et la
page 38 s'en sert.

\emph{Conventions.} Pour un foncteur $f : X \to Y$ et un objet $b$ de $Y$, on
note $X_{b}$ la fibre, $X_{/b}$ la catégorie des couples $(x, f(x) \to b)$ et
${}_{b\backslash}X$ celle des couples $(x, b \to f(x))$\footnote{Ce sont ses
graphies, celles de la page 8. Aujourd'hui on écrit souvent $f/b$ et $b\backslash
f$, ou $f \downarrow b$ et $b \downarrow f$.}. $X^{\circ}$ est la catégorie
opposée, $e$ la catégorie ponctuelle (il écrit aussi $\Delta_{0}$), $\Delta^{1}$
la catégorie $\{0 \to 1\}$. Une \emph{fibration} est une catégorie fibrée au
sens de Grothendieck (relèvements cartésiens), une \emph{cofibration} une
catégorie cofibrée (relèvements cocartésiens) ; ce ne sont pas les fibrations
et cofibrations d'une structure de modèles. Une sous-catégorie pleine
$U \subset X$ est \emph{ouverte} si toute flèche $x \to u$ de but dans $U$ a sa
source dans $U$ (un crible), \emph{fermée} si toute flèche de source dans $U$ a
son but dans $U$ (un cocrible) ; le complémentaire d'un ouvert est fermé.

$\mathcal{W}$ désigne la classe des équivalences faibles de (Cat), qu'il
soumet aux axiomes « Loc »\footnote{Il souligne la lettre $W$ à partir de la
page 14, et non aux pages 2 à 8 ; on lit partout la même classe. Aux pages 21
à 40 il omet parfois le soulignement au milieu d'une formule où le contexte
désigne la même classe.} ; une catégorie $X$ est \emph{asphérique} si
$X \to e$ est dans $\mathcal{W}$, et un foncteur $f$ est asphérique si toutes
les $X_{/b}$ le sont. Pour une petite catégorie $A$, $\widehat{A}$ est la
catégorie des préfaisceaux d'ensembles sur $A$\footnote{Il écrit tantôt
$A^{\wedge}$, tantôt $\widehat{A}$.}, $i_{A} : \widehat{A} \to (\mathrm{Cat})$
le foncteur $X \mapsto A_{/X}$, et $\mathcal{W}_{A} = i_{A}^{-1}(\mathcal{W})$.
On pose $\mathrm{Hot} = \mathrm{Hot}(\mathcal{W}) = \mathcal{W}^{-1}(\mathrm{Cat})$,
$\mathrm{Hot}_{A} = \mathcal{W}_{A}^{-1}\widehat{A}$, et $\gamma$ désigne le
foncteur de localisation.

\emph{Deux homonymies à ne pas confondre.} La lettre $\mathrm{C}$ désigne la
classe des cofibrations, et $C(f)$, $C(Y)$, en italique, le cylindre d'un
foncteur et le cône d'une catégorie (page 10). La lettre $\mathrm{S}$ désigne
la classe des foncteurs « de Serre », sans rapport avec un ensemble $S$.

\pagerange{2}{8}
\section*{Les classes de foncteurs (pages 2 à 8)}

\pagerange{2}{4}
\subsection*{Les treillis (pages 2 et 4)}

La page 2 dessine sans un mot trois treillis de classes de foncteurs ; la page
4 en reprend les deux premiers sous leur forme la plus complète, entourés de
remarques. Les classes de base sont
$\mathrm{F}$ (fibrations), $\mathrm{C}$ (cofibrations), $\mathrm{Sm}$ (foncteurs
lisses), $\mathrm{Pr}$ (foncteurs propres), $\mathrm{S}$ (« Serre map, i.e.
Hot-fibration ») et $\mathrm{B} = \mathrm{F} \cap \mathrm{C}$ (bifibrations) ;
leurs définitions sont à la page 8, ci-dessous. Viennent ensuite trois classes
plus faibles, $\mathrm{WSF}$, $\mathrm{WSC}$, $\mathrm{WSB}$, que la page 2
appelait $\mathrm{Spr}F$, $\mathrm{Spr}C$, $\mathrm{Spr}B$\footnote{Les sigles
ne sont développés nulle part. La page 4 glose $\mathrm{WSF}$ par « Serre
\emph{fibrations} » et $\mathrm{WSC}$ par « Serre precofibrations » (lecture
incertaine), chacun précédé d'un mot biffé ; « weakly Serre » est une
conjecture de lecture, non un mot de la page.}. Rangées par inclusion, la
flèche allant de la plus petite à la plus grande :

\begin{tikzcd}[column sep=tiny, row sep=small, nodes={font=\scriptsize}]
 & & & & \mathrm{B} \arrow[dl] \arrow[dr] & & & & \\
 & & & \mathrm{F}\cap\mathrm{Pr} \arrow[dl] \arrow[dr] & & \mathrm{Sm}\cap\mathrm{C} \arrow[dl] \arrow[dr] & & & \\
 & & \mathrm{F}\cap\mathrm{S} \arrow[dl] \arrow[dr] & & \mathrm{Sm}\cap\mathrm{Pr} \arrow[dl] \arrow[dr] & & \mathrm{S}\cap\mathrm{C} \arrow[dl] \arrow[dr] & & \\
 & \mathrm{F} \arrow[dr] & & \mathrm{Sm}\cap\mathrm{S} \arrow[dl] \arrow[dr] \arrow[dddd, dashed] & & \mathrm{S}\cap\mathrm{Pr} \arrow[dl] \arrow[dr] \arrow[dddd, dashed] & & \mathrm{C} \arrow[dl] & \\
 & & \mathrm{Sm} & & \mathrm{S} \arrow[d] & & \mathrm{Pr} & & \\
 & & & & \mathrm{WSB} \arrow[d] & & & & \\
 & & & & \mathrm{WSF}\cap\mathrm{WSC} \arrow[dl] \arrow[dr] & & & & \\
 & & & \mathrm{WSF} \arrow[dr] & & \mathrm{WSC} \arrow[dl] & & & \\
 & & & & \mathrm{Fl}(\mathrm{Cat}) & & & &
\end{tikzcd}

Les deux longues flèches, de $\mathrm{Sm} \cap \mathrm{S}$ vers $\mathrm{WSF}$
et de $\mathrm{S} \cap \mathrm{Pr}$ vers $\mathrm{WSC}$, sont hachurées sur la
page et tracées ici en tirets\footnote{La page porte aussi, au crayon, des
étiquettes sur quatre flèches : $\mathrm{F} \cap \mathrm{WSC}$,
$\mathrm{Sm} \cap \mathrm{WSC}$, $\mathrm{WSF} \cap \mathrm{C}$,
$\mathrm{WSF} \cap \mathrm{Pr}$, et le mot « str » sur
$\mathrm{F} \cap \mathrm{S} \to \mathrm{F}$. Les cercles dont il entoure
$\mathrm{F}$, $\mathrm{C}$, $\mathrm{Sm}$, $\mathrm{S}$, $\mathrm{Pr}$ ne
sont pas reproduits. À la page 2, la version antérieure place, au centre du
premier losange, un $F' \cap C'$ qui n'étiquette aucune flèche.}. Les
inclusions qu'il encadre sont
\[
\mathrm{F} \subsetneq \mathrm{Sm}, \quad \mathrm{C} \subsetneq \mathrm{Pr},
\quad \mathrm{Sm} \cap \mathrm{Pr} \subsetneq \mathrm{S}, \quad
\mathrm{WSC} \cap \mathrm{Pr} \subset \mathrm{S}, \quad
\mathrm{WSF} \cap \mathrm{Sm} \subset \mathrm{S},
\]
d'où $\mathrm{WSC} \cap \mathrm{Pr} = \mathrm{S} \cap \mathrm{Pr}$ et
$\mathrm{WSF} \cap \mathrm{Sm} = \mathrm{S} \cap \mathrm{Sm}$, puisque
$\mathrm{S} \subset \mathrm{WSC} \cap \mathrm{WSF}$. L'inclusion
$\mathrm{Sm} \cap \mathrm{Pr} \subset \mathrm{S}$ est l'analogue du théorème
d'Ehresmann : un morphisme lisse et propre est localement trivial. Il note
aussi, avec un point d'exclamation, $\mathrm{WSB} \neq \mathrm{WSF} \cap
\mathrm{WSC}$, et que $\mathrm{B}$ n'est contenue ni dans $\mathrm{WSF}$ ni
dans $\mathrm{WSC}$ (ni a fortiori $\mathrm{C}$ dans $\mathrm{WSF}$,
$\mathrm{F}$ dans $\mathrm{WSC}$).

Pour la stabilité par changement de base, les classes du premier étage
(jusqu'à $\mathrm{S}$) le sont par changement de base quelconque, les carrés
du treillis étant cartésiens\footnote{« stable under any base change ; squares
cartesian », dans une accolade à droite du treillis. Une seconde note marginale,
« stable by intersection except $WSF\cap C$ and $WSC\cap Pr$, and dually
$F\cap WSC$, $Sm\cap WSC$ and $WSF\cap WSC$ », dit quels sommets ne sont pas
l'intersection de leurs deux majorants ; on la rapporte sans la développer.} ;
$\mathrm{WSF}$ l'est par changement de base propre, $\mathrm{WSC}$ par
changement de base lisse.

Le second treillis est l'intersection du premier avec $\mathcal{W}$. On y lit
$\mathcal{W} \cap \mathrm{F} \cap \mathrm{S} = \mathrm{T} \cap \mathrm{F}$,
$\mathcal{W} \cap \mathrm{S} \cap \mathrm{C} = \mathrm{T} \cap \mathrm{C}$, et
ainsi de suite, où $\mathrm{T}$ est la classe des foncteurs
\emph{universellement} dans $\mathcal{W}$ (tout changement de base est dans
$\mathcal{W}$) ; puis, sous $\mathcal{W} \cap \mathrm{S} = \mathrm{T} \cap
\mathrm{S}$\footnote{La page prolonge l'égalité par un troisième membre, lu
avec doute $UW$ ou $\mathcal{U}W$, qui n'est défini nulle part.},
\[
\mathrm{Tot\,As} = \mathcal{W} \cap \mathrm{WSB}, \qquad
\mathrm{As}^{\circ} = \mathcal{W} \cap \mathrm{WSF}, \qquad
\mathrm{As} = \mathcal{W} \cap \mathrm{WSC},
\]
au-dessus de $\mathcal{W}$ lui-même, qui est stable par changement de base le
long d'un foncteur de Serre --- c'est la définition de $\mathrm{S}$ --- tandis
que $\mathrm{As}$ l'est par changement de base lisse et $\mathrm{As}^{\circ}$
par changement de base propre.

Ces égalités sont à lire avec précaution. $\mathrm{As}$ est la classe des
foncteurs asphériques, ceux dont toutes les $X_{/b}$ sont asphériques, et
$\mathrm{As}^{\circ}$ celle des foncteurs coasphériques. Un foncteur asphérique
est dans $\mathcal{W}$ (c'est le théorème A de Quillen, que l'axiome Loc 3
généralise) et dans $\mathrm{WSC}$, ses $X_{/b}$ étant toutes asphériques ;
on a donc $\mathrm{As} \subset \mathcal{W} \cap \mathrm{WSC}$. L'inclusion
inverse dit que les $X_{/b}$ d'un foncteur de $\mathcal{W} \cap \mathrm{WSC}$
sont asphériques : c'est un énoncé de type théorème B de
Quillen\footnote{L'égalité est celle de la page. Pour l'équivalence faible
usuelle elle est vraie : sous l'hypothèse de $\mathrm{WSC}$, le théorème B
identifie $X_{/b}$ à la fibre homotopique de $f$, contractile si $f$ est une
équivalence faible. Pour un localisateur fondamental quelconque, le dossier
n'établit pas de théorème B, et l'on ne retient que l'inclusion.}. De même,
la page affirme qu'un foncteur de Serre est dans $\mathcal{W}$ si et seulement
si ses fibres sont asphériques, si et seulement s'il est universellement dans
$\mathcal{W}$ ; le sens « fibres asphériques $\Rightarrow$ $\mathcal{W}$ » est
une conséquence du théorème A une fois les $X_{/b}$ comparées aux fibres, le
sens réciproque est encore du type théorème B.

Restent des questions, qu'il cerne d'un trait : $\mathrm{WSC}$ et
$\mathrm{WSB}$ sont-elles stables par changement de base le long d'une
cofibration de Serre, même « triviale » ? Il précise ce qu'il cherche : un
$f \in \mathcal{W} \cap \mathrm{WSB}$ (« totalement asphérique ») et un
$g \in \mathrm{T} \cap \mathrm{C}$ (cofibration à fibres triviales) dont le
changement de base $f'$ ne soit pas dans $\mathrm{WSB}$. Et encore : même à
fibres asphériques, la condition de Serre ne peut-elle se tester sur les seuls
changements de base $\Delta_{1} \to X$ ? Aucune n'est tranchée dans le
dossier.

Le troisième treillis de la page 2, au crayon, est le même que le premier
intersecté avec une classe $\mathrm{strh}$, que la page 7 glose « strict
homotopism »\footnote{À la page 7, au milieu d'une page blanche : « $TF$ »
biffé, une flèche verticale $X \to Y$ et les mots « strict homotopism »
soulignés d'un trait ondulé. Dans ce troisième treillis, $\mathrm{Sm}$ et
$\mathrm{Pr}$ deviennent $\mathrm{str}Sm$, $\mathrm{str}Pr$ ; aucune de ces
variantes strictes n'est définie.}.

\pagerange{7}{8}
\subsection*{Les caractérisations (page 8)}

Sur une feuille de papier kraft, en français, chaque rubrique met une classe à
gauche et sa duale à droite. Soit $f : X \to Y$.

\emph{Cofibrations.} Pour une flèche $b \to b'$ de $Y$ et un objet $a$ de la
fibre $X_{b}$, considérons la catégorie des relèvements de $b \to b'$ de source
$a$ : ses objets sont les flèches $a \to x$ de $X$ au-dessus de $b \to b'$.
Elle a un objet initial pour tout $(b \to b', a)$ exactement quand $f$ est
précofibré --- toute flèche de la base a un relèvement cocartésien de source
donnée ---, et $f$ est une cofibration si de plus ces relèvements se
composent\footnote{La page donne trois formes équivalentes : $X_{/b} \to
X_{/b'}$ « cofinal à gauche », ${}_{b'\backslash}X \to {}_{b\backslash}X$ à
fibres non vides ayant un objet initial, et l'énoncé sur les relèvements. Ni
la stabilité des relèvements par composition ni le passage de l'une à l'autre
ne sont écrits ; on retient la forme sur les relèvements, la plus directe, en
y ajoutant la condition de composition. La lettre de la rubrique, repassée, se
lit aussi bien $E$ que $C$ ; on lit $C$ parce que la colonne de droite porte
$F$.}. Dualement pour les fibrations, la colonne de droite restant vide.

\emph{Foncteurs propres et lisses.} La définition est obtenue en affaiblissant
« objet initial » en « asphérique » : $f$ est \emph{propre} si, pour toute flèche
$b \to b'$ de $Y$ et tout $a \in X_{b}$, la catégorie des relèvements de
$b \to b'$ de source $a$ est asphérique. La page en donne des formes
équivalentes : le foncteur d'inclusion $X_{b} \to X_{/b}$ est coasphérique
(dans $\mathrm{As}^{\circ}$) ; les foncteurs ${}_{a\backslash}X \to
{}_{b\backslash}Y$, pour $b = f(a)$, sont à fibres asphériques, ou sont de
Serre ; et enfin les images directes supérieures $R^{i}f_{*}$ commutent à tout
changement de base\footnote{Les abréviations serrées « Rif », « Rig » sont lues
avec doute $R^{i}f_{*}$, $R^{i}g_{*}$ ; la lettre de la rubrique, surchargée,
est lue $Pr$ parce que la colonne de droite porte $Sm$, et le mot qui précède
« coasph. », « strongly », est lui aussi incertain. Une note cernée en marge
ajoute que ${}_{a\backslash}X \to {}_{b\backslash}Y$ est propre, peut-être « à
fibres discrètes ».}. L'inclusion $\mathrm{C} \subset \mathrm{Pr}$ se voit
sur la première forme : pour une cofibration, $X_{b} \to X_{/b}$ a un adjoint à
gauche, et un adjoint à droite est coasphérique. Dualement, $f$ est
\emph{lisse} si $f^{\circ}$ est propre, et le changement de base par un
foncteur lisse commute aux $R^{i}g_{*}$. Les noms viennent de l'analogie avec
les théorèmes de changement de base propre et lisse en cohomologie
étale\footnote{En dessous, en anglais : « [$\Longleftarrow$ ? base changes by
$f$ preserves cocartesian squares] », une implication qu'il se demande et ne
tranche pas.}.

\emph{Foncteurs de Serre.} $f \in \mathrm{S}$ si, pour tout foncteur
$g : Y'' \to Y'$ au-dessus de $Y$ qui est dans $\mathcal{W}$, le changement de
base $g \times_{Y} X : X'' \to X'$ est encore dans $\mathcal{W}$ : changer de
base le long de $f$ préserve les équivalences faibles au-dessus de
$Y$\footnote{C'est, presque mot pour mot, la définition que la théorie des
modèles donne aujourd'hui d'un morphisme « sharp » (Rezk, 1998) ; le nom est
nôtre et postérieur. Dans une catégorie de modèles propre à droite, les
fibrations sont sharp, comme ici les fibrations « universellement » bien
comportées.}. On a $\mathrm{S} = \mathrm{S}^{\circ}$, $\mathrm{C} =
\mathrm{F}^{\circ}$, $\mathrm{Sm} = \mathrm{Pr}^{\circ}$, et
$\mathrm{T} \cap \mathrm{Pr} \subset \mathrm{S}$,
$\mathrm{T} \cap \mathrm{Sm} \subset \mathrm{S}$.

\emph{Les classes faibles.} $f \in \mathrm{WSC}$ si, pour toute flèche
$b \to b'$ de $Y$, le foncteur $X_{/b} \to X_{/b'}$ est dans
$\mathcal{W}$ : c'est l'hypothèse du théorème B de Quillen, sous l'une de ses
deux formes duales. La page en donne deux autres formes : la condition de
Serre restreinte à une classe de bases $Y'' \to Y'$ qu'elle dit « fibrés »
au-dessus de $Y$\footnote{« fibrés » est écrit au-dessus d'un mot biffé et lu
sans certitude.}, et la condition que la cofibration associée
$\mathrm{Cofib}(X/Y) \to Y$ (page 10) soit de Serre. $\mathrm{WSF}$ est la
notion duale, avec ${}_{b'\backslash}X \to {}_{b\backslash}X$ et
$\mathrm{Fib}(X/Y)$. Ces classes sont l'analogue des quasi-fibrations de Dold
et Thom : un foncteur qui, remplacé par une (co)fibration, devient de
Serre\footnote{L'analogie est la nôtre. Une note marginale, à hauteur de la
dernière forme, se lit avec doute « $R^{i}f_{*}$ (c$^{\mathrm{t}}$)
loc. const. ». La rubrique ajoute que $\mathrm{WSC}$ est stable par changement
de base « propres lisses » (lecture incertaine, sur « fibrations » biffé),
« mais non par $g \in S$ ? ». Enfin la page encadre « $WSC \cap Pr = S$ » ; on
lit l'inclusion $\mathrm{WSC} \cap \mathrm{Pr} \subset \mathrm{S}$, comme à la
page 4 : l'égalité avec $\mathrm{S}$ tout entière est fausse, $\mathrm{S}$
n'étant pas contenue dans $\mathrm{Pr}$.}.

$\mathrm{WSB}$ est la variante à deux côtés : pour une flèche $u : b_{0} \to
b_{1}$ de $Y$, soit $X[u] = {}_{b_{0}\backslash}X_{/b_{1}}$ la catégorie des
factorisations de $u$ à travers $f$ ; $f \in \mathrm{WSB}$ si, pour tout
composé $u' : b'_{0} \to b_{0} \xrightarrow{u} b_{1} \to b'_{1}$, le foncteur
$X[u] \to X[u']$ est dans $\mathcal{W}$ --- en d'autres termes, si les
équivalences faibles ${}_{b_{0}\backslash}Y_{/b_{1}} \to
{}_{b'_{0}\backslash}Y_{/b'_{1}}$ le restent après produit fibré avec
$X$\footnote{La page continue par deux équivalences, coupées par la
déchirure du bas de la feuille : la condition sur les ${}_{b'\backslash}X \to
{}_{b\backslash}X$ « sur $Y$ », et sur les $X_{/b} \to X_{/b'}$, « coloc. sur
$Y$ » (lecture incertaine).}.

\pagerange{10}{13}
\section*{Factorisations, immersions, cylindres (pages 10 à 13)}

\subsection*{Les deux factorisations par les catégories comma}

Soit $f : X \to Y$. Notons $\mathrm{Fib}(f)$ la catégorie des triplets
$(b, x, u : b \to f(x))$, avec les flèches évidentes. Elle porte trois
foncteurs :
\[
\begin{array}{lll}
f^{\flat} : \mathrm{Fib}(f) \to Y, & (b, x, u) \mapsto b, & \text{une fibration,} \\
\mu_{f} : X \to \mathrm{Fib}(f), & x \mapsto (f(x), x, \mathrm{id}), & \text{pleinement fidèle,} \\
\rho_{f} : \mathrm{Fib}(f) \to X, & (b, x, u) \mapsto x, & \text{une rétraction,}
\end{array}
\]
avec $f = f^{\flat}\mu_{f}$ et $\rho_{f}\mu_{f} = \mathrm{id}_{X}$. Le premier
est une fibration, dont la fibre en $b$ est ${}_{b\backslash}X$ --- « but not
yet Serre » ; $\mu_{f}$ est pleinement fidèle ; et $\rho_{f}$ est adjoint à
gauche de $\mu_{f}$, puisque les flèches $(b, x', u) \to \mu_{f}(x)$
correspondent aux flèches $x' \to x$. C'est la factorisation d'une application
par son espace des chemins : $f$ s'écrit comme une équivalence d'homotopie
suivie d'une fibration\footnote{$f^{\flat}$ est ce qu'on appelle aujourd'hui
la fibration libre engendrée par $f$ (Street, 1974) ; le nom est le nôtre. La
page dit « $(\mu_{f}, \rho_{f})$ adjoint » sans préciser de quel côté ; le
calcul donne $\rho_{f} \dashv \mu_{f}$. L'exposant de $f$, lu $\flat$, est
l'un des signes incertains du dossier.}. Elle ajoute que $\rho_{f}$ est une
cofibration de Serre « triviale », à fibres les $Y_{/f(x)}$, qui ont un objet
final, et un homotopisme ; la factorisation duale passe par
$\mathrm{Cofib}(f)$, catégorie des $(b, x, f(x) \to b)$, avec une cofibration
$\mathrm{Cofib}(f) \to Y$ à fibres les $X_{/b}$.

\subsection*{Immersions et classes discrètes}

Les pages 10, 11 et 13 disposent les classes en croix --- fibrations et
cofibrations sur les bras, bifibrations et Serre en haut, monomorphismes et
classes discrètes en bas, lisses et propres aux extrémités --- et identifient
les immersions\footnote{La page 11 est une mise au propre de la page 10, que
la page 13 reprend en tableau. À la page 10, « monoe » (sic) pour
« monos ».} :
\begin{itemize}
\item une immersion ouverte $Y \hookrightarrow X$ est la même chose qu'un
foncteur $\varphi : X \to \Delta^{1}$, avec $Y = \varphi^{-1}(0)$, ou encore
qu'un crible de $X$ ; c'est une fibration discrète qui est un monomorphisme ;
\item une immersion fermée est la même chose qu'un foncteur
$\varphi : X \to \Delta^{1}$ avec $Y = \varphi^{-1}(1)$, un cocrible ; c'est
une cofibration discrète monomorphe\footnote{Pour l'immersion ouverte la page
écrit $\varphi : X^{\circ} \to \Delta^{1}$ et $Y = \varphi^{-1}(\{1\})$, ce
qui revient au même par $\Delta_{1}^{\circ} \simeq \Delta_{1}$ ; elle note
aussi $\Delta^{1} = [\emptyset \to \{\emptyset\}]$, l'ensemble ordonné des
parties de l'ensemble vide.} ;
\item une fibration discrète au-dessus de $X$ est un préfaisceau
$X^{\circ} \to (\mathrm{Ens})$, une cofibration discrète un foncteur
$X \to (\mathrm{Ens})$.
\end{itemize}
Au bas de la page 11, il note que « discret » et Serre donnent les revêtements
étales, que les immersions à la fois ouvertes et fermées sont les inclusions de
sommandes directes (de réunions de composantes connexes), et qu'un revêtement
de Serre trivial est un isomorphisme\footnote{La première ligne se lit
« discrete $\cap$ Serre ($=$ étale covering) $=$ smooth $\cap$ proper
$\cap$ Serre », un mot biffé au milieu ; les premiers mots des deux lignes
suivantes sont surchargés, et la dernière se lit peut-être « discrete $\cap$
trivial Serre $=$ iso ».} ; puis qu'un foncteur propre et discret est une
cofibration, un foncteur lisse et discret une fibration, mais qu'un foncteur
propre de Serre n'est pas nécessairement une cofibration, ni lisse de Serre
une fibration.

\subsection*{Cylindres et cônes}

Le \emph{cylindre} $C(f)$ de $f : Y \to X$ est la catégorie dont les objets
sont ceux de $Y$ et ceux de $X$, avec les flèches de $Y$, celles de $X$, et
$\mathrm{Hom}(y, x) = \mathrm{Hom}_{X}(f(y), x)$, aucune flèche n'allant de $X$
vers $Y$\footnote{La construction n'est pas décrite sur la page, qui n'en donne
que les propriétés ; on la reconstitue d'après elles. C'est ce qu'on appelle
aussi le \emph{collage} de $f$.}. On a une factorisation
\[
Y \xrightarrow{\ f^{\natural}\ } C(f) \xrightarrow{\ \lambda_{f}\ } X ,
\qquad \sigma_{f} : X \to C(f), \quad \lambda_{f}\sigma_{f} = \mathrm{id}_{X},
\]
où $f^{\natural}$ est une immersion ouverte, $\sigma_{f}$ une section qui est
une immersion fermée, et $\lambda_{f}$ un épimorphisme adjoint à gauche de
$\sigma_{f}$ ; ses fibres ont un objet final, et $\sigma_{f}$ est un
homotopisme. La page 13 demande si $(\lambda_{f}, \sigma_{f})$ est une
adjonction : c'est le cas\footnote{$\mathrm{Hom}_{X}(\lambda_{f}(c), x) =
\mathrm{Hom}_{C(f)}(c, \sigma_{f}(x))$ pour $c$ objet de $Y$ ou de $X$ ; la
vérification est la nôtre. Le tableau de la page 13 classe aussi
« $f^{\sharp} \in$ fibration $\cap$ mono » et « $\sigma_{f} \in$
cofibration $\cap$ mono » ; l'exposant lu $\sharp$ pourrait être un
$\natural$, et la case décrit l'immersion ouverte $f^{\natural}$, qui est bien
une fibration monomorphe : on lit $f^{\natural}$. L'exposant $\natural$ est
lui-même incertain.}. Toute flèche se factorise ainsi en une immersion ouverte
suivie d'une flèche de $\mathcal{W}$ ; dualement, $C'(f)$ donne une immersion
fermée suivie d'une flèche à section ouverte, avec
$C'(f)^{\circ} = C(f^{\circ})$.

Le \emph{cône} $C(Y)$ est $C(Y \to e)$ : on adjoint à $Y$ un objet final strict,
et $Y$ y est ouverte ; $C'(Y)$ adjoint un objet initial strict, et $Y$ y est
fermée.

\pagerange{14}{15}
\section*{Le programme (pages 14 et 15)}

« Après l'histoire d'asphéricité », c'est-à-dire une fois acquis les axiomes
Loc 1 à Loc 3, Loc 1 étant pris sous forme faible --- les isomorphismes sont
dans $\mathcal{W}$ ; si deux des trois flèches $f$, $g$, $gf$ sont dans
$\mathcal{W}$, la troisième aussi ; si $gf = \mathrm{id}$ et $fg \in
\mathcal{W}$, alors $f, g \in \mathcal{W}$ ---, la page dresse un programme en
six rubriques\footnote{Les axiomes Loc 2 et Loc 3 ne sont pas rappelés ici ;
la page 25 nomme Loc 3 « condition de localisation » et s'en sert comme d'un
théorème A relatif : un foncteur $X \to Y$ au-dessus de $J$ est dans
$\mathcal{W}$ dès que les $X_{/j} \to Y_{/j}$ le sont. Ce sont, dans le
vocabulaire actuel, les axiomes d'un \emph{localisateur fondamental}, notion
qui apparaît dans \emph{Pursuing Stacks} et que Cisinski a étudiée ensuite.} :
\begin{itemize}
\item[(0)] mettre en œuvre les axiomes Loc, jusqu'aux fibrations et
cofibrations\footnote{Lecture incertaine ; la ligne est reliée par une flèche
à « … transitions fib et cofib. », dont le premier mot ne se lit pas.} ;
\item[(1)] la dualité, y compris $\mathcal{W} = \mathcal{W}^{\circ}$ ;
\item[(2)] la commutation aux fibrations, cofibrations, immersions ouvertes et
fermées, cylindres, cônes ;
\item[(3)] le « sorite des $\int$ » d'un diagramme $X_{1} \xleftarrow{f_{1}} Y
\xrightarrow{f_{0}} X_{0}$ ;
\item[(4)] le théorème « de la somme amalgamée », dans (Cat) et dans les
$\widehat{A}$ ;
\item[(5)] l'application à $\mathrm{Hot}(\mathcal{W})$ et aux
$\mathrm{Hot}_{A}$, pour $A$ catégorie test locale : sommes et produits ;
\item[(6)] la saturation forte de $\mathcal{W}$ et des
$\mathcal{W}_{A}$\footnote{« forte » est ajouté au-dessus de la ligne. Les
numéros sont cerclés et plusieurs sont repassés ($3$ sur $4$, $4$ sur $3$,
$5$ sur un autre chiffre) ; la moitié inférieure de la page 14 est très
raturée.}.
\end{itemize}

Le « sorite des $\int$ » porte sur la construction de Grothendieck d'un
diagramme $X_{1} \xleftarrow{f_{1}} Y \xrightarrow{f_{0}} X_{0}$, vu comme un
foncteur de la catégorie $\cdot \leftarrow \cdot \rightarrow \cdot$ vers
(Cat) : la catégorie $\int(f_{0}, f_{1})$ dont les objets sont ceux de $Y$, de
$X_{0}$ et de $X_{1}$, avec les flèches de chacune et
$\mathrm{Hom}(y, x_{k}) = \mathrm{Hom}_{X_{k}}(f_{k}(y), x_{k})$. Elle s'envoie
sur la somme amalgamée $X_{0} \amalg_{Y} X_{1}$, et c'est un modèle du
recollement homotopique\footnote{Que la construction de Grothendieck calcule
la colimite homotopique est le théorème de Thomason (1979) ; la rubrique ne le
cite pas, et nous ne disons pas qu'il le connaissait.}. La rubrique en énonce
la fonctorialité --- si les trois composantes d'un morphisme de diagrammes sont
dans $\mathcal{W}$, le foncteur induit entre les $\int$ l'est ---, le cas où
$f_{0}$ est un isomorphisme, où les flèches canoniques sont des homotopismes,
et le cas où $f_{0} \in \mathcal{W}$, où elles sont dans $\mathcal{W}$. Elle
esquisse ensuite l'énoncé central : si $f_{0}$, $f_{1}$ sont des cofibrations
(resp. des fibrations) et les fibres convenablement asphériques, le foncteur
$\int(f_{0}, f_{1}) \to X_{0} \amalg_{Y} X_{1}$ est une cofibration (resp. une
fibration) à fibres asphériques, donc est dans $\mathcal{W}$ ; d'où, si
$f_{0} \in \mathcal{W}$, la flèche opposée du carré cocartésien est dans
$\mathcal{W}$\footnote{La page marque le carré de la rubrique (4) « cart » ;
l'énoncé porte sur la somme amalgamée, et dans la situation des pages 18 à 21
le carré est à la fois cartésien et cocartésien. On le lit cocartésien. Le
passage de la page 14 qui précède la rubrique (4) est en grande partie
illisible, et l'énoncé n'en est qu'esquissé ; les pages 18 à 20 le reprennent
et le démontrent.}.

Une \emph{question} suit, pour un carré cocartésien
$X' = X \amalg_{Y} Y'$ le long d'une immersion ouverte ou fermée
$i : Y \hookrightarrow X$ : si $f : Y \to Y'$ est une cofibration, la flèche
$g : X \to X'$ l'est-elle, et si $f \in \mathcal{W}$, a-t-on $g \in
\mathcal{W}$ ? « Il est possible qu'on n'ait pas $g \in \mathcal{W}$ [mais
j'ignore si $i \in \mathcal{W}$ implique que $i' \in \mathcal{W}$] ». La page
22 répond par la négative. Dans un $\widehat{A}$, au contraire, la rubrique
annonce le théorème pour un carré cocartésien le long d'un monomorphisme $i$ :
$i \in \mathcal{W}_{A} \Rightarrow i' \in \mathcal{W}_{A}$ et
$f \in \mathcal{W}_{A}$ pour le côté opposé à un $g \in
\mathcal{W}_{A}$\footnote{La source de la première ligne du carré est écrite
$Y'$, pour $Y$ sans doute.}. La page 15 ajoute l'application aux sommes et
produits dans $\mathrm{Hot}(\mathcal{W})$ et dans les $\mathrm{Hot}_{A}$, et
la saturation forte.

\pagerange{16}{16}
\section*{Le sorite des recollements (page 16)}

Soient $X$ une catégorie, $U$ un ouvert et $Y$ le fermé complémentaire.
Aucune flèche ne va de $Y$ vers $U$, et $X$ est déterminée par $U$, $Y$ et le
profoncteur
\[
H_{X} : U^{\circ} \times Y \to (\mathrm{Ens}), \qquad H_{X}(u, y) =
\mathrm{Hom}_{X}(u, y),
\]
contravariant en $u$ et covariant en $y$ : $X$ est le collage de
$H_{X}$\footnote{« Collage » et « profoncteur » (ou « distributeur », chez
Bénabou) sont des noms modernes ; la page n'en emploie aucun. Le cylindre
$C(f)$ de la page 10 est le collage du profoncteur $\mathrm{Hom}_{X}(f(-),
-)$.}. Le sorite décrit tout en ces termes :
\begin{itemize}
\item[a)] un foncteur $X \to Z$ est la donnée de $f_{U} : U \to Z$, de
$f_{Y} : Y \to Z$ et d'une transformation $H_{X}(u, y) \to
\mathrm{Hom}_{Z}(f_{U}(u), f_{Y}(y))$, naturelle en $u$ et $y$ ;
\item[b)] un foncteur $f : X' \to X$ compatible aux décompositions
($X' = U' \cup Y'$, de profoncteur $H_{X'}$) est la donnée de
$f_{U'} : U' \to U$, de $f_{Y'} : Y' \to Y$ et d'une transformation naturelle
$H_{X'}(u', y') \to H_{X}(f_{U'}(u'), f_{Y'}(y'))$ ;
\item[c)] pour une catégorie $T$, $T \times X$ se décompose en
$T \times U$ et $T \times Y$, avec $H_{T \times X}((t, u), (t', y)) \simeq
\mathrm{Hom}_{T}(t, t') \times H_{X}(u, y)$ ;
\item[d)] une homotopie $h : \Delta^{1} \times X' \to X$ est compatible aux
décompositions si $h^{-1}(U) = \Delta^{1} \times U'$ et $h^{-1}(Y) =
\Delta^{1} \times Y'$, avec $U' = h_{i}^{-1}(U)$, $Y' = h_{i}^{-1}(Y)$ pour
$i = 0, 1$ ; elle est alors donnée par $h_{U'} : \Delta^{1} \times U' \to U$,
$h_{Y'} : \Delta^{1} \times Y' \to Y$ et une application
$\mathrm{Hom}(s, t) \times H_{X'}(u', y') \to H_{X}(h_{U'}(s, u'), h_{Y'}(t,
y'))$, naturelle en $(s, u')$ et $(t, y')$\footnote{La page écrit
$h^{-1}(Y')$ pour $h^{-1}(Y)$ et, deux fois, $\Delta$ pour $\Delta^{1}$ ; les
égalités $U' = h_{i}^{-1}(U)$, $Y' = h_{i}^{-1}(Y)$ rendent deux guillemets de
répétition.}.
\end{itemize}

\pagerange{18}{20}
\section*{Sommes amalgamées le long d'une immersion (pages 18 à 20)}

\subsection*{Quand le carré est-il cocartésien ?}

Soient $X'$ une catégorie, $Y'$ un fermé, $U'$ l'ouvert complémentaire,
$f : X \to X'$ un foncteur, $Y = f^{-1}(Y')$ et $U = f^{-1}(U')$, et supposons
que $f$ induise un isomorphisme $U \simeq U'$. Le carré
\[
\begin{array}{ccc}
Y & \xrightarrow{\ i\ } & X \\
{\scriptstyle g}\downarrow & & \downarrow{\scriptstyle f} \\
Y' & \xrightarrow{\ i'\ } & X'
\end{array}
\]
n'est pas cocartésien en général : il suffit de prendre $X = U'$ et $f$
l'inclusion, d'où $Y = \emptyset$. Il l'est si $f$ est une cofibration. Les
deux membres ont en effet les mêmes objets et les mêmes flèches à l'intérieur
de $U'$ et de $Y'$, et il s'agit de voir que, pour $u' \in U'$ (relevé en
$u \in U$) et $y' \in Y'$,
\[
\varinjlim_{(y,\, g(y) \to y') \in Y_{/y'}} \mathrm{Hom}_{X}(u, y)
\longrightarrow \mathrm{Hom}_{X'}(u', y')
\]
est bijectif. Comme $g$, changement de base de $f$, est une cofibration, la
fibre $Y_{y'}$ est cofinale dans $Y_{/y'}$ et la colimite se calcule sur
elle ; or, à $u$ et à une flèche $u' \to y'$ fixés, la catégorie des
relèvements a un objet initial, et elle est donc connexe et non vide. La page
dit : « Oui, car (i) surjectif et (ii) pour $u \to y'$ fixé
… »\footnote{La phrase s'interrompt là ; l'argument est complété ici.
Dans la formule, la page indexe la colimite par $Y_{/y'}$ et, par une
accolade, par la fibre $Y_{y'}$.}.

Se donner une telle cofibration $f$, dont les fibres au-dessus de $U'$ sont
ponctuelles, revient à se donner un foncteur $X' \to (\mathrm{Cat})$ de
restriction constante de valeur $e$ sur $U'$, c'est-à-dire un foncteur
$\Phi : Y' \to (\mathrm{Cat})$, qui définit $g : Y \to Y'$, et des
applications $\mathrm{Hom}(u', y') \to \mathrm{Ob}\,\Phi(y')$, naturelles en
$u'$ et $y'$.

\subsection*{Le théorème}

Il voudrait alors que $\int(i, g) \to X' = X \amalg_{Y} Y'$ soit dans
$\mathcal{W}$, et même asphérique\footnote{« même » est rendu sur la page par
le signe $=$ ; le passage est marqué d'un double trait en marge.}. Ces
conditions sont stables par tout changement de base sur $X'$, et il se demande
si le changement de base commute à $\int$ --- « Ce diagramme est-il
cartésien ? À vérifier ! » ---, ce qui ferait de $\int(i, g) \to X'$ un
foncteur universellement dans $\mathcal{W}$. La question reste ouverte, mais
on peut s'en passer : on vérifie directement que $\int(i, g) \to X'$ est une
cofibration, de fibre $e$ au-dessus de $u' \in U'$ et, au-dessus de
$y' \in Y'$, isomorphe au cône $C(Y_{y'})$, qui a un objet final. D'où
l'énoncé de la page 20.

\textbf{Théorème.} Soit, dans (Cat), un carré cocartésien et cartésien
\[
\begin{array}{ccc}
Y & \xrightarrow{\ i\ } & X \\
{\scriptstyle g}\downarrow & & \downarrow{\scriptstyle f} \\
Y' & \xrightarrow{\ i'\ } & X'
\end{array}
\]
où $i'$ et $f$ sont des cofibrations, et supposons que pour tout objet $x'$ de
$X'$, l'une des deux fibres $f^{-1}(x')$, $i'^{-1}(x')$ soit équivalente à la
catégorie ponctuelle. Alors le foncteur canonique $\int(i, g) \to X'$ est une
cofibration à fibres asphériques ; il est donc dans $\mathcal{W}$.

Les fibres sont les $\int$ des diagrammes $Y'_{x'} \leftarrow Y'_{x'} \times
X_{x'} \rightarrow X_{x'}$ formés des deux projections\footnote{L'égalité
$Y_{x'} = Y'_{x'} \times X_{x'}$, que la page écrit (surchargée) dans
l'intégrale, demande que le carré soit cartésien ; l'énoncé de la page ne le
dit pas. L'hypothèse est satisfaite dans la situation des pages 18 et 19, où
$Y = f^{-1}(Y')$, et dans celle de la page 21, que la page déclare cartésienne.
On l'ajoute à l'énoncé.}. Si l'un des deux facteurs est ponctuel, à
équivalence près, une telle fibre est de la forme $\int(A = A, A \to e)$, ou
de la forme symétrique. Cette catégorie est $C(A) \amalg_{A} (A \times
\Delta^{1})$ : un cône sur $A$ prolongé par un cylindre, dont le cône est un
rétracte par déformation ; elle est donc asphérique, « on gagne ! »\footnote{Une
note marginale ajoute qu'au lieu de « ponctuelle » il suffit
qu'une fibre soit asphérique ; elle n'est pas démontrée. Un passage
suivant, enclos d'un trait et barré, voulait voir directement que
$\int(A \to A)$ et $\int(A \to e)$ sont asphériques « par propriété des
recollements » ; il est trop fragmentaire pour être lu. Un « Lemme », écrit
en oblique au pied de la page, est inachevé après ses premiers mots.}

\textbf{Corollaire.} \emph{Sous les mêmes hypothèses, $i \in \mathcal{W}
\Rightarrow i' \in \mathcal{W}$, et $g \in \mathcal{W} \Rightarrow f \in
\mathcal{W}$.}

La page énonce le corollaire sans démonstration. Il découle du théorème et de
la fonctorialité des $\int$ (page 14) : si $i \in \mathcal{W}$, le morphisme
de diagrammes de $(Y = Y \xrightarrow{g} Y')$ vers $(Y \xrightarrow{i} X,\ Y \xrightarrow{g} Y')$
 induit une flèche de $\mathcal{W}$ entre leurs $\int$,
et $Y' \to \int(\mathrm{id}_{Y}, g)$ est un homotopisme ; la composée
$Y' \to \int(i, g) \to X'$ est $i'$, et les deux facteurs sont dans
$\mathcal{W}$. De même pour $f$\footnote{La déduction est la nôtre.}.

\pagerange{21}{22}
\section*{Dans les préfaisceaux, et une mise en garde (pages 21 et 22)}

\subsection*{Le théorème dans un $\widehat{A}$}

Soit, dans $\widehat{A}$, un carré cocartésien
\[
\begin{array}{ccc}
Y & \hookrightarrow & X \\
\downarrow & & \downarrow \\
Y' & \hookrightarrow & X'
\end{array}
\]
où $Y \to X$ est un monomorphisme ; il est alors aussi cartésien, et
$Y' \to X'$ est un monomorphisme\footnote{La page ajoute dans l'interligne,
soulignés, « et aussi cartésien » ; dans un topos de préfaisceaux, une somme
amalgamée le long d'un monomorphisme est toujours un carré cartésien, et le
monomorphisme se transporte. La remarque est nôtre.}. Le foncteur $i_{A}$
transforme ce carré en un carré de (Cat) qui est encore cocartésien, puisque
$i_{A}$ commute aux limites inductives --- il a un adjoint à droite --- et
cartésien, puisqu'il commute aux produits fibrés. Les flèches
$A_{/X} \to A_{/X'}$, $A_{/Y} \to A_{/Y'}$ sont des fibrations discrètes, et
$A_{/Y} \to A_{/X}$, $A_{/Y'} \to A_{/X'}$, fibrations discrètes monomorphes,
sont des immersions ouvertes. On est dans la situation du théorème de la page
20, sous sa forme duale, celle où $i'$ et $f$ sont des fibrations : au-dessus
d'un objet de $A_{/X'}$ hors de $A_{/Y'}$, la fibre de $f$ est un point, et
au-dessus d'un objet de $A_{/Y'}$, la fibre de $i'$ en est un\footnote{La
vérification des hypothèses sur les fibres est la nôtre ; elle tient à ce
que $X'(a) = X(a) \amalg_{Y(a)} Y'(a)$ pour tout objet $a$ de $A$. La page dit
seulement « On est donc dans la situation déjà énoncée », et la fin de sa
phrase, un $\int$ suivi de $A_{/X'}$, est d'une lecture incertaine.}. D'où :

\textbf{Corollaire} (en marge : « Thm »). \emph{Pour un tel carré, si
$Y \to X$ est dans $\mathcal{W}_{A}$, $Y' \to X'$ l'est ; si $Y \to Y'$ est
dans $\mathcal{W}_{A}$, $X \to X'$ l'est.}

\subsection*{Mais attention}

Dans (Cat), il est faux qu'une somme amalgamée le long d'une immersion fermée
$i$ transporte les équivalences faibles --- ni $i \in \mathcal{W} \Rightarrow
i' \in \mathcal{W}$, ni $g \in \mathcal{W} \Rightarrow f \in \mathcal{W}$ ---,
et cela même pour l'équivalence faible usuelle, même si $g$ est une
cofibration et $Y'$ ponctuelle. L'exemple, pour une catégorie $Z$, empile deux
carrés cocartésiens :
\[
\begin{array}{ccc}
Z & \hookrightarrow & C'_{2}(Z) \\
\downarrow & & \downarrow \\
C'_{1}(Z) & \hookrightarrow & \Sigma(Z) \\
{\scriptstyle g}\downarrow & & \downarrow{\scriptstyle f} \\
e & \hookrightarrow & \Theta
\end{array}
\]
où $C'_{1}(Z)$ et $C'_{2}(Z)$ sont deux exemplaires du cône à sommet initial,
dans lesquels $Z$ est fermée, $\Sigma(Z)$ leur recollement le long de $Z$ ---
une suspension ---, et $\Theta$ la catégorie obtenue en écrasant $C'_{1}(Z)$
en un point\footnote{La page écrit $C(Z)$, $\check{C}(Z)$ et $\Theta(Y)$, et
marque les flèches de $Z$ dans $C(Z)$ et de $C(Z)$ dans $\Sigma(Z)$
« immersion fermée » ; pour qu'elles le soient, les deux cônes doivent avoir
leur sommet initial, le $C'$ de la page 10, et on le lit ainsi. Les indices
des sommets, lus $u'_{0}$, $w_{0}$, $y'_{0}$, sont incertains.}. Le
foncteur $g : C'_{1}(Z) \to e$ est dans $\mathcal{W}$, le cône étant
asphérique. La catégorie $\Theta$ a deux objets, le sommet du second cône et
le point, et l'ensemble des flèches du premier vers le second est $\pi_{0}(Z)$ ;
elle est donc équivalente à $\Delta^{1}$, et asphérique, quand $Z$ est
connexe. Dans ce cas $f \in \mathcal{W}$ si et seulement si $\Sigma(Z)$ est
asphérique, ce qui n'a pas toujours lieu : pour $Z$ un groupe $G$, vu comme
groupoïde à un objet, le nerf de $\Sigma(Z)$ a le type de la suspension de
$BG$, qui n'est pas contractile dès que l'abélianisé de $G$ n'est pas nul ---
par exemple pour $G = \mathbf{Z}$, où l'on obtient la sphère
$S^{2}$\footnote{La page dit : « prendre p.ex. $Z$ un groupoïde réduit à un
objet, avec $\pi_{1} \neq 1$ ». La condition $\pi_{1} \neq 1$ ne suffit pas :
pour un groupe acyclique non trivial, comme celui de Higman, $BG$ a l'homologie
d'un point et sa suspension, simplement connexe et acyclique, est
contractile. Il faut une suspension non contractile, ce qu'assure par exemple
$H_{1}(G) \neq 0$. L'exemple réfute $g \in \mathcal{W} \Rightarrow f \in
\mathcal{W}$ ; la page annonce aussi l'échec de $i \in \mathcal{W} \Rightarrow
i' \in \mathcal{W}$ sans en donner d'exemple.}.

Une note marginale avance qu'« il semble vrai » que si $g$ est une
1-équivalence --- une équivalence sur les groupoïdes fondamentaux ---, alors
$f$ l'est aussi\footnote{Note écrite en diagonale, lue en partie avec doute ;
elle n'est pas reprise. Dans l'exemple ci-dessus, $f$ est bien une
1-équivalence, $\Sigma(Z)$ étant simplement connexe pour $Z$ connexe. La
littérature ultérieure contourne l'obstacle en restreignant les inclusions le
long desquelles on recolle : ce sont les « Dwyer maps » de la structure de
modèles de Thomason (1980), dont les sommes amalgamées sont homotopiques ; la
page n'en parle pas.}.

\pagerange{23}{25}
\section*{La catégorie homotopique comme catégorie de fractions (pages 23 à 25)}

Soit $A$ une catégorie test locale\footnote{L'hypothèse est en marge,
séparée d'une seconde note : « itou dans Cat ». La page 25 dit pourquoi elle
est nécessaire.}. Notons $\overline{\widehat{A}}$ la catégorie localisée de
$\widehat{A}$ par les homotopismes relatifs à $\mathcal{W}_{A}$, et
$\overline{\mathcal{W}_{A}}$ l'image de $\mathcal{W}_{A}$ ; alors
$\mathrm{Hot}_{A} \simeq \overline{\mathcal{W}_{A}}^{\,-1}\,
\overline{\widehat{A}}$.

\textbf{Corollaire.} \emph{$\overline{\mathcal{W}_{A}}$ admet un calcul de
fractions dans lequel toute flèche de $\mathrm{Hot}_{A}$ s'écrit
$\gamma(t)^{-1}\gamma(u)$, pour un diagramme $X \xrightarrow{u} Y'
\xleftarrow{t} Y$ de $\overline{\widehat{A}}$ avec $t \in
\overline{\mathcal{W}_{A}}$.}\footnote{La page dit « calcul des fractions à
gauche » ; la page 26, pour la même situation, dit « à droite ». Les deux
noms s'échangent d'un auteur à l'autre ; on désigne le calcul par sa forme,
celle que les deux arguments utilisent.}

Il y a deux conditions à vérifier. La première : un diagramme
$Y' \xleftarrow{g} Y \xrightarrow{i} X$ avec $g \in
\overline{\mathcal{W}_{A}}$ se complète en un carré commutatif
$f \circ i = i' \circ g$, avec $f \in \overline{\mathcal{W}_{A}}$. On factorise
$i$ en $Y \xrightarrow{i_{1}} X_{1} \xrightarrow{h} X$, avec $i_{1}$
monomorphe et $h$ homotopisme, ce qui ramène au cas où $i$ est un
monomorphisme, et l'on prend pour $X'$ la somme amalgamée ; le corollaire de la
page 21 donne $f \in \mathcal{W}_{A}$.

La seconde est la simplification : si $f, g : X \rightrightarrows Y$ et
$s : X' \to X$ dans $\overline{\mathcal{W}_{A}}$ vérifient $fs = gs$ dans
$\overline{\widehat{A}}$, il existe $t : Y \to Y'$ dans
$\overline{\mathcal{W}_{A}}$ avec $tf = tg$. On peut supposer $s$ monomorphe.
L'égalité $fs = gs$ est donnée par une homotopie $I \times X' \to Y$. Soit
$M$ la somme amalgamée de $I \times X'$ et de $X \sqcup X$ au-dessus de
$X' \sqcup X'$. La flèche $X \sqcup X \to M$ est dans $\mathcal{W}_{A}$, comme
somme amalgamée de $X' \sqcup X' \to X \sqcup X$ le long du monomorphisme
$X' \sqcup X' \to I \times X'$ ; donc $M \to I \times X$, monomorphe, est dans
$\mathcal{W}_{A}$ ; et l'on pousse $M \to Y$ le long de $M \to I \times X$ pour
obtenir $t : Y \to Y'$, dans $\mathcal{W}_{A}$, et une homotopie
$I \times X \to Y'$ entre $tf$ et $tg$\footnote{La page donne le diagramme, en
tirets pour les flèches à construire, et non le texte ; le récit est le
nôtre. La lettre de $Y \to Y'$ est lue $t$, et pourrait être un $r$.}.

L'argument n'utilise que quatre propriétés d'une classe $\mathfrak{M}$ de
flèches d'une catégorie $M$ munie de $\mathcal{W}$ (ici
$M = \widehat{A}$ et $\mathfrak{M}$ les monomorphismes) :
\begin{itemize}
\item[(i)] toute flèche se factorise en une flèche de $\mathfrak{M}$ suivie
d'un homotopisme ;
\item[(ii)] dans un carré cocartésien le long d'une flèche de $\mathfrak{M}$,
le côté opposé à une flèche de $\mathcal{W}$ est dans $\mathcal{W}$ ;
\item[(iii)] il existe une famille génératrice d'intervalles $I_{\alpha}$ avec
$X \sqcup X \to X \times I_{\alpha}$ dans $\mathfrak{M}$ pour tout $X$ --- des
intervalles « séparants » ;
\item[(iv)] les sommes finies existent, et $f, g \in \mathcal{W} \Rightarrow
f \sqcup g \in \mathcal{W}$.
\end{itemize}
Une note marginale l'énonce en général : pour $(M, \mathcal{W})$ muni d'un tel
$\mathfrak{M}$, les homotopismes étant dans $\mathcal{W}$, on a
$\mathcal{W}^{-1}M = \overline{\mathcal{W}}^{\,-1}\overline{M}$ et
$\overline{\mathcal{W}}$ admet ce calcul de fractions\footnote{C'est, à la
dualité près, le schéma des catégories d'objets fibrants de K. Brown (1973),
où la catégorie homotopique s'obtient par fractions à partir de la catégorie
des classes d'homotopie ; le rapprochement est le nôtre. Une autre note
marginale dit les conditions « inspirées par le cas de (Cat) et non celui de
$A^{\wedge}$, où (iii) est inutile » ; la note se termine sur des mots
illisibles.}. Il se demande d'ailleurs si l'argument vaut dans (Cat) ; la
page 22 interdit d'y prendre pour $\mathfrak{M}$ les immersions fermées, et le
dossier ne dit pas quelle classe conviendrait.

La page 25 dit ce qu'a coûté l'argument. Il a fallu un intervalle séparant $I$
de $\widehat{A}$ définissant une structure homotopique compatible avec
$\mathcal{W}_{A}$, et pour cela que $I \to e$ soit dans
$\mathcal{W}_{A}$, ou du moins que les deux sections $\delta_{0}$,
$\delta_{1} : e \to I$ aient même image dans $\mathrm{Hot}_{A}$ --- « donc,
pratiquement, que $A$ soit une catégorie-test \emph{locale} »\footnote{Les
deux notes marginales de la page 25, l'une écrite presque verticalement, sont
en grande partie illisibles ; un mot biffé en haut à droite ne se lit pas
avec certitude.}. Il a fallu aussi la stabilité de $\mathcal{W}_{A}$ par
sommes finies, propriété (iv) : elle se ramène à (Cat), puisque $i_{A}$
commute aux sommes, et là elle résulte de Loc 3 appliqué à un foncteur
$X \to Y$ au-dessus de l'ensemble d'indices $J$, vu comme catégorie
discrète.

\pagerange{26}{30}
\section*{Les sommes dans $\mathrm{Hot}$ (pages 26 à 30)}

\textbf{Corollaire 1.} \emph{Les sommes existent dans $\mathrm{Hot}
= \mathrm{Hot}(\mathcal{W})$, et les foncteurs canoniques $(\mathrm{Cat}) \to
\mathrm{Hot}$ et $\widehat{A} \to \mathrm{Hot}_{A} \to \mathrm{Hot}$ y
commutent.}

L'énoncé est démontré pour les sommes \emph{finies} ; pour les sommes
quelconques, que la page semble viser, il reste ouvert dans le
dossier\footnote{Le mot qui qualifie les sommes est lu avec doute
« arbitraires », et une note marginale en oblique, à moitié illisible, porte
« \emph{finies} inutile » et « quelconque ». Les pages 28 à 30 montrent que
le cas infini n'est pas acquis.}.

\emph{Le principe.} Pour une catégorie $M$ et une classe $\mathcal{W}$
admettant le calcul de fractions de la page 23, on a
$\mathrm{Hom}_{\mathcal{W}^{-1}M}(X, Y) = \varinjlim \mathrm{Hom}_{M}(X, Y')$,
la colimite, filtrante, portant sur les flèches $Y \to Y'$ de $\mathcal{W}$ ;
le foncteur $M \to \mathcal{W}^{-1}M$ commute donc aux limites inductives
finies qui existent dans $M$, et en particulier aux sommes finies\footnote{La
page dit « commute aux \emph{lim} finies » ; pour ce calcul de fractions, ce
sont les limites inductives finies, et c'est à elles, les sommes, qu'elle
l'applique.}. On l'applique à $(\overline{\mathrm{Cat}},
\overline{\mathcal{W}})$, où $\overline{\mathrm{Cat}}$ est la catégorie des
classes d'homotopie de foncteurs ; il faut que les sommes finies y existent et
que $(\mathrm{Cat}) \to \overline{\mathrm{Cat}}$ y commute. C'est le

\textbf{Lemme} (page 27). \emph{Soit $(M, h)$ une structure homotopique par
intervalles, avec $M$ stable par sommes finies et les sommes finies
universelles. Alors $M \to \overline{M}$ commute aux sommes finies.}

Il s'agit de voir que deux flèches $f = (f_{j})$, $g = (g_{j}) :
\coprod_{J} X_{j} \rightrightarrows Y$, $J$ fini, sont homotopes dès que
chaque $f_{j}$ l'est à $g_{j}$. Les homotopies élémentaires entre $f_{j}$ et
$g_{j}$ se mettent bout à bout en une homotopie paramétrée par un même
intervalle composé $I_{N}$, $N$ assez grand pour tous les $j$ --- c'est là que
sert la finitude de $J$ ---, et comme $I_{N} \times \coprod X_{j} \simeq
\coprod I_{N} \times X_{j}$, elles se recollent en une homotopie entre $f$ et
$g$.

\emph{Le cas infini.} Le même énoncé est faux pour les sommes infinies :
$(\mathrm{Cat}) \to \overline{\mathrm{Cat}}$ n'y commute pas. Soit
$Y = I_{\infty}$ le zigzag infini $\bullet \to \bullet \leftarrow \bullet \to
\bullet \leftarrow \cdots$, réunion d'une infinité d'intervalles
élémentaires\footnote{La page dit « somme amalgamée d'une infinité de copies
de $X_{n}$ », où $X_{n} = \Delta_{0}$ ; ce sont des copies de l'intervalle
élémentaire qu'on recolle. La parenthèse ouverte avant l'exemple n'est pas
refermée.}, $X_{n} = e$ pour $n \in \mathbf{N}$, et $f_{n}$, $g_{n} : e \to
Y$ les points $0$ et $2n$. Chaque $f_{n}$ est homotope à $g_{n}$, mais au
moyen d'une chaîne de longueur croissante avec $n$ ; une homotopie de
$\coprod X_{n}$ à valeurs dans $Y$ n'a qu'une longueur finie, et $f$, $g$
ne sont donc pas homotopes. Ainsi $\coprod X_{n}$ n'est pas une somme dans
$\overline{\mathrm{Cat}}$. « Je suppose pourtant que dans
$\mathrm{Hot}(\mathcal{W})$, les sommes quelconques existent et que
$\mathrm{Cat} \to \mathrm{Hot}(\mathcal{W})$ y commute ?? »

Cela reviendrait à la bijectivité de
\[
\mathrm{Hom}_{\mathrm{Hot}}\Bigl(\coprod_{J} X_{j}, Y\Bigr) \longrightarrow
\prod_{J} \mathrm{Hom}_{\mathrm{Hot}}(X_{j}, Y),
\]
et, « avec le calcul des fractions, ni l'injectivité ni la surjectivité n'est
claire ». Pour la surjectivité, la difficulté est de trouver une même flèche
$Y \to Y'$ de $\mathcal{W}$ qui représente à la fois toutes les classes
$\varphi_{j}$ ; chacune en fournit une, $Y \to Y'_{j}$, et il propose de les
réunir par sommes amalgamées au-dessus de $Y$ sur les parties finies
$J_{0} \subset J$, puis par passage à la limite, ce qui donne encore une flèche
de $\mathcal{W}$ --- « Donc il y a quand même surjectivité ! » L'argument
repose sur un

\textbf{Lemme} (page 30). \emph{Soit $(\alpha_{j} : Y \to Y'_{j})_{j \in J}$
une famille d'équivalences faibles. Il existe une équivalence faible
$\alpha : Y \to Y'$, immersion ouverte ou fermée si les $\alpha_{j}$ le sont,
et des $\beta_{j} : Y'_{j} \to Y'$ avec $\beta_{j}\alpha_{j} = \alpha$ ; si les
$\alpha_{j}$ sont des immersions ouvertes ou fermées, on peut prendre pour $Y'$
la somme amalgamée des $Y'_{j}$ au-dessus de $Y$.}

Il n'est qu'esquissé\footnote{Un schéma en marge factorise $\alpha_{j}$ par une
immersion $Y \hookrightarrow \overline{Y}'_{j}$ suivie de
$i_{j} : \overline{Y}'_{j} \to Y'_{j}$, avec une flèche $s_{j}$ en sens
inverse et la relation $s_{j}i_{j} = \mathrm{id}_{\overline{Y}'_{j}}$, sans
dire d'où vient la factorisation. Avec le cylindre des pages 10 et 11, elle
existe, mais la relation est $i_{j}s_{j} = \mathrm{id}_{Y'_{j}}$, et l'égalité
$\beta_{j}\alpha_{j} = \alpha$ ne vaut qu'à homotopie près. Que la somme
amalgamée, finie puis infinie, de ces immersions soit encore dans
$\mathcal{W}$ demanderait dans (Cat) les hypothèses de fibration de la page
20, que la page 22 montre nécessaires, puis le passage à la réunion des pages
38 et 39 ; ni l'un ni l'autre n'est invoqué ici.}. L'injectivité, enfin, se
ramène à ceci : si $f, g : \coprod X_{j} \rightrightarrows Y$ ont des
composantes $f_{j}$, $g_{j}$ homotopes, a-t-on $\gamma(f) = \gamma(g)$ ? « Mais
attention on n'a pas en général $f \sim g$ ! » Il faudrait une condition « de
type Kan » sur $Y$, ou sur un $Y'$ avec $Y \to Y'$ dans $\mathcal{W}$ ; le
dossier s'arrête là.

\pagerange{31}{36}
\section*{Les produits finis (pages 31 à 36)}

\textbf{Corollaire 2.} \emph{Les produits finis existent dans $\mathrm{Hot}$,
et les foncteurs canoniques $(\mathrm{Cat}) \to \mathrm{Hot}$ et
$\widehat{A} \to \mathrm{Hot}$ y commutent, ce dernier pour $A$ catégorie test
stricte, ou plus généralement totalement asphérique.}\footnote{« il suffit
même que $A$ soit tot. $\mathcal{W}$-asph. », en marge. L'abréviation lue
« tot. », ici et jusqu'à la page 35, ressemble à « LA » ou « tA » ; elle est
lue « totalement » d'après la condition de la page 33. Une catégorie $A$ est
\emph{totalement asphérique} si elle est asphérique et si les produits
$a \times b$ de deux représentables sont asphériques dans $\widehat{A}$.}

Pour $(\mathrm{Cat})$ le point est « clair si on traite directement », un
produit fini de flèches de $\mathcal{W}$ étant dans $\mathcal{W}$ (page 35).
Pour $\widehat{A}$, on est ramené à voir que $\widehat{A} \to \mathrm{Hot}_{A}$
commute aux produits finis : que
\[
\mathrm{Hom}_{\mathrm{Hot}_{A}}\Bigl(X, \prod Y_{j}\Bigr) \longrightarrow
\prod \mathrm{Hom}_{\mathrm{Hot}_{A}}(X, Y_{j})
\]
est bijectif pour toute famille finie $(Y_{j})$. \emph{Surjectivité} : un
élément du second membre est donné par des $f_{j} : X \to Y'_{j}$, avec
$Y_{j} \to Y'_{j}$ dans $\mathcal{W}_{A}$ ; $f = (f_{j}) : X \to \prod Y'_{j}$
définit un antécédent, \emph{si l'on sait} que $\prod Y_{j} \to \prod Y'_{j}$
est dans $\mathcal{W}_{A}$. \emph{Injectivité} : soient
$f, g : X \rightrightarrows Y'$, $\alpha : Y = \prod Y_{j} \to Y'$ dans
$\mathcal{W}_{A}$, et pour chaque $j$ des $\beta_{j} : Y_{j} \to Y'_{j}$ dans
$\mathcal{W}_{A}$ et $\varphi_{j} : Y' \to Y'_{j}$ avec
$\beta_{j}\,\mathrm{pr}_{j} = \varphi_{j}\alpha$ et $\varphi_{j}f$ homotope à
$\varphi_{j}g$. Avec $\varphi = (\varphi_{j})$ et $\beta = \prod \beta_{j}$,
on a $\varphi\alpha = \beta$ et $\varphi f$ homotope à $\varphi g$ ; si
$\beta$ est dans $\mathcal{W}_{A}$, $\gamma(\varphi)$ est inversible et
$\gamma(f) = \gamma(g)$. Dans les deux sens, tout repose donc sur la
stabilité de $\mathcal{W}_{A}$ par produits finis ; et cela marcherait pour
les produits infinis si l'on savait qu'un produit infini de flèches de
$\mathcal{W}_{A}$ est dans $\mathcal{W}_{A}$.

\subsection*{Le lemme des pages 32 à 35}

La page 32 énonce, avec un point d'interrogation, un lemme de « conditions
équivalentes » pour une petite catégorie $A$ :
\begin{itemize}
\item[(i)] $A$ est totalement asphérique ;
\item[(ii)] $\mathcal{W}_{A}$ est stable par produits finis ;
\item[(ii bis)] pour $f : X \to Y$ dans $\mathcal{W}_{A}$ et $Z$ dans
$\widehat{A}$, $f \times \mathrm{id}_{Z}$ est dans $\mathcal{W}_{A}$ ;
\item[(iii)] $\widehat{A} \to \mathrm{Hot}_{A}$ commute aux produits finis ;
\item[(iv)] $\widehat{A} \to \mathrm{Hot}$ commute aux produits binaires ;
\item[(iv bis)] $\widehat{A} \to \mathrm{Hot}$ commute aux produits finis.
\end{itemize}
Elles ne sont pas toutes équivalentes, et les pages 33 à 35 le montrent
elles-mêmes. Ce qu'elles établissent est ceci\footnote{À la page 32 le lemme
s'intitule « Conditions équiv. ($\Rightarrow$) », et le schéma de la page 34
groupe (i), (iv), (iv bis) d'un côté, (ii), (ii bis), (iii) de l'autre, avec
des doubles flèches. Sa propre démonstration n'obtient que les énoncés
ci-dessous, et la page 35 montre que (iv) n'implique pas (i). La page 32 écrit
(iv) avec $\mathrm{Hot}_{A}(\mathcal{W})$, la page 34 avec
$\mathrm{Hot}(\mathcal{W})$ ; la démonstration porte sur
$\mathrm{Hot}$.} :
\begin{itemize}
\item (ii) $\Leftrightarrow$ (ii bis), et il suffit de prendre pour $Z$ les
représentables $a$ : la condition se teste localement sur $Z$ ;
\item (ii) $\Rightarrow$ (iii), par l'argument ci-dessus ; la réciproque vaut
quand $\mathcal{W}_{A}$ est fortement saturé, une flèche inversible dans
$\mathrm{Hot}_{A}$ étant alors dans $\mathcal{W}_{A}$ ;
\item (i) $\Rightarrow$ (ii bis) : dans le carré formé de
$X \times a \to Y \times a$ et de $X \to Y$, les projections
$X \times a \to X$, $Y \times a \to Y$ sont dans $\mathcal{W}_{A}$ quand $A$
est totalement asphérique ; inversement, si $A$ est asphérique, (ii bis)
appliqué à $a \to e$ donne (i). Donc (i) $\Leftrightarrow$ (ii) $+$ « $A$
asphérique » ;
\item (iv) équivaut à l'asphéricité des $a \times b$ : comme
$(\mathrm{Cat}) \to \mathrm{Hot}$ commute déjà aux produits finis, (iv)
signifie que $A_{/X \times Y} \to A_{/X} \times A_{/Y}$ est dans $\mathcal{W}$
pour tous $X$, $Y$ ; pour $X = a$, $Y = b$ c'est l'asphéricité de $a \times b$ ;
et inversement, pour un objet $\alpha = (a \to X, b \to Y)$ du second membre,
$(A_{/X \times Y})_{/\alpha} \simeq A_{/a \times b}$, si bien que le foncteur
est asphérique, donc dans $\mathcal{W}$ ;
\item (iv bis) $\Leftrightarrow$ (i), le produit vide demandant que
$A = A_{/e}$ soit asphérique ; et (i) $\Leftrightarrow$ (iv) $+$ « $A$
asphérique ».
\end{itemize}
La page 35 précise que l'implication (i) $\Rightarrow$ (iv) est stricte :
(iv) est stable par sommes de catégories, (i) ne l'est pas, une somme n'étant
pas connexe. Une catégorie connexe qui satisfait (ii) est-elle asphérique ?
\emph{Non} : pour $A$ un groupoïde connexe dont les automorphismes forment un
groupe $G \neq 1$, et $\mathcal{W} = \mathcal{W}_{0}$ l'équivalence faible
usuelle, $\mathcal{W}_{A}$ est formé des seuls isomorphismes --- un morphisme
de $G$-ensembles qui induit une bijection sur les orbites et des isomorphismes
sur les stabilisateurs est un isomorphisme --- et elle est stable par produits
finis, alors que $A$, de nerf $BG$, n'est pas asphérique\footnote{La raison
donnée pour « $\mathcal{W}_{A}$ formé des seuls iso » est la nôtre ; la page
renvoie à l'équivalence faible « habituelle ».}.

Les produits infinis restent en suspens. Il faudrait savoir qu'un produit de
flèches de $\mathcal{W}$ est dans $\mathcal{W}$, ou le même énoncé dans un
$\widehat{A}$, le second impliquant le premier ; « déjà pour l'équiv. faible
habituelle $\mathcal{W}_{0}$, ce n'est pas évident ». Il suggère de passer
par l'asphéricité d'un produit infini de catégories asphériques, et juge la
question « moins importante de toutes façons que la question des sommes
infinies ».

\pagerange{38}{39}
\section*{Filtrations (pages 38 et 39)}

\textbf{Théorème.} \emph{Soient $I$ un ensemble ordonné dans lequel toute
partie non vide a une borne inférieure, et $X$, $Y$ deux catégories munies de
familles croissantes $(X_{i})_{i \in I}$, $(Y_{i})_{i \in I}$ de
sous-catégories ouvertes (resp. fermées), avec $X = \bigcup X_{i}$,
$Y = \bigcup Y_{i}$, et $X_{\inf j_{\alpha}} = \bigcap X_{j_{\alpha}}$, de même
pour $Y$. Soit $f : X \to Y$ un foncteur tel que $f(X_{i}) \subset Y_{i}$, et
notons $f_{i} : X_{i} \to Y_{i}$ ses restrictions. Si les $f_{i}$ sont dans
$\mathcal{W}$, $f$ est dans $\mathcal{W}$.}\footnote{L'énoncé de la page
suppose seulement $I$ ordonné (« filtrant » est biffé) ; la note marginale
dit qu'« il suffit » que les bornes inférieures existent et que la
filtration les respecte, et la démonstration se place d'emblée sous cette
hypothèse. On l'a mise dans l'énoncé. « ouvertes » est lu avec doute.}

Le cas fermé se ramène au cas ouvert par passage aux catégories
opposées\footnote{Cela suppose $\mathcal{W} = \mathcal{W}^{\circ}$, la
rubrique (1) du programme de la page 14, qui n'est pas établie dans le
dossier. Pour l'équivalence faible usuelle, le nerf de $X^{\circ}$ a la même
réalisation que celui de $X$, et le point est acquis.}. Dans le cas ouvert, la
filtration de $X$ est la même chose qu'un foncteur $\varphi : X \to I$ :
\[
\varphi(x) = \inf I(x), \qquad I(x) = \{\, i \in I \mid x \in X_{i} \,\},
\]
et $\varphi(x)$ est le plus petit élément de $I(x)$, qui est l'ensemble des
$j \geq \varphi(x)$. Pour une flèche $x \to y$, on a $I(y) \subset I(x)$,
puisque les $X_{i}$ sont des cribles, donc $\varphi(x) \leq \varphi(y)$ ; et
$X_{i} = \varphi^{-1}(I_{\leq i})$, c'est-à-dire $x \in X_{i}$ si et seulement
si $\varphi(x) \leq i$. Il note qu'il faudra expliciter la correspondance
réciproque : pour un $I$ où les bornes inférieures des parties non vides
existent, les foncteurs $X \to I$ correspondent aux filtrations croissantes
$i \mapsto X_{i}$ par des ouverts qui commutent à ces bornes et dont $X$ est
la réunion --- ce qui équivaut à $\varinjlim X_{i} \simeq X$.

Comme $I$ est un ensemble ordonné, la catégorie $X_{/i}$ relative à $\varphi$
est exactement $X_{i}$, et de même $Y_{/i} = Y_{i}$ pour le foncteur
$\psi : Y \to I$ de la filtration de $Y$ ; les $f_{i}$ sont donc les foncteurs
$f_{/i} : X_{/i} \to Y_{/i}$, et Loc 3 conclut\footnote{La page dit « par
L 3 ». L'hypothèse $f(X_{i}) \subset Y_{i}$ ne donne que $\psi \circ f \leq
\varphi$ : le triangle formé de $f$, $\varphi$ et $\psi$ ne commute qu'à une
transformation naturelle $\psi f \Rightarrow \varphi$ près, et c'est bien
cette transformation qui définit les $f_{/i}$. Loc 3 sert donc sous une forme
admettant un tel triangle, que la page ne distingue pas de la forme stricte.
Si $f^{-1}(Y_{i}) = X_{i}$ pour tout $i$, le triangle commute et la forme
stricte suffit.}.

\textbf{Corollaire.} \emph{De même dans un $\widehat{A}$, pour
$\mathcal{W}_{A}$}, puisque $i_{A}$ commute aux limites inductives
quelconques et aux produits fibrés.

\emph{« Cor ? »} Si les $X_{i}$ sont asphériques et $I$ connexe, $X$ est
asphérique : on prend $Y_{i} = e$ pour tout $i$\footnote{Le point
d'interrogation est le sien. Sous l'hypothèse du théorème, $I$, partie non
vide de lui-même, a un plus petit élément, ce qui le rend connexe, et
l'énoncé découle alors du théorème. La remarque est nôtre.}.

\pagerange{39}{40}
\section*{La saturation (pages 39 et 40)}

\textbf{Théorème.} \emph{$\mathcal{W}$ est saturé : un foncteur $f$ tel que
$\gamma(f)$ soit inversible dans $\mathrm{Hot}$ est dans $\mathcal{W}$.}

\textbf{Corollaire.} \emph{Pour toute $A$, $\mathcal{W}_{A}$ est saturé.}

C'est la rubrique (6) du programme, la saturation « forte »\footnote{Il dit
seulement « saturé » à la page 39 et « saturation forte » à la page 15 ; le
sens est celui de la seconde. Le corollaire est énoncé sans démonstration ; il
découle du théorème en une ligne : si $\gamma_{A}(f)$ est inversible dans
$\mathrm{Hot}_{A}$, son image par le foncteur $\mathrm{Hot}_{A} \to
\mathrm{Hot}$ induit par $i_{A}$ l'est, donc $i_{A}(f) \in \mathcal{W}$,
c'est-à-dire $f \in \mathcal{W}_{A}$. La démonstration est la nôtre.}. Une
première démonstration, sur la page 39, est barrée de deux longs traits : elle
partait de $f : X \to Y$ et de $g : Y \to X'$, $f' : X' \to Y'$, avec
$\gamma(gf) = \gamma(\alpha)$ et $\gamma(f'g) = \gamma(\beta)$ pour
$\alpha$, $\beta$ dans $\mathcal{W}$, et cherchait à se ramener à
$gf = \alpha$, $f'g = \beta$ et à des immersions ouvertes.

La page 40 donne la démonstration qui aboutit. Soit $f_{0} : X_{0} \to X_{1}$
avec $\gamma(f_{0})$ inversible. Par le cylindre de la page 10, on peut
supposer $f_{0}$ immersion ouverte. Par le calcul des fractions, il existe
$f_{1} : X_{1} \to X_{2}$ avec $f_{1}f_{0} \in \mathcal{W}$ ; alors
$\gamma(f_{1})$ est inversible, et l'on peut supposer $f_{1}$ immersion
ouverte\footnote{Le calcul de fractions invoqué est celui de
$(\overline{\mathrm{Cat}}, \overline{\mathcal{W}})$, que la marge de la page
23 étend à (Cat) d'un « itou » et que le dossier n'établit que dans les
$\widehat{A}$.}. De proche en proche, on obtient une suite d'immersions
ouvertes
\[
X_{0} \xrightarrow{\ f_{0}\ } X_{1} \xrightarrow{\ f_{1}\ } X_{2}
\xrightarrow{\ f_{2}\ } \cdots
\]
dont les composés de deux flèches consécutives sont dans $\mathcal{W}$. Soit
$X_{\infty} = \varinjlim_{n} X_{n}$, et $i_{n} : X_{n} \to X_{\infty}$ ; on a
$i_{1}f_{0} = i_{0}$, et il suffit de prouver que $i_{0}$ et $i_{1}$ sont dans
$\mathcal{W}$. Or $X_{\infty}$ est aussi la limite de la sous-suite
$X_{0} \to X_{2} \to X_{4} \to \cdots$, dont les flèches sont dans
$\mathcal{W}$, et $i_{0}$ est la limite des $X_{0} \to X_{2n}$, qui sont dans
$\mathcal{W}$ : « donc on gagne !! » ; de même pour $i_{1}$, en partant de
$f_{1}$. Le passage à la limite est le théorème des filtrations des pages 38
et 39, appliqué à $I = \mathbf{N}$, à la filtration constante de $X_{0}$ et à
la filtration de $X_{\infty}$ par les ouverts $X_{2n}$\footnote{La page ne
dit pas ce qui justifie le passage à la limite ; le renvoi au théorème des
pages 38 et 39 est le nôtre. Ici le triangle est strict, $X_{0}$ étant
envoyé dans chaque $X_{2n}$, et la forme stricte de Loc 3 suffit. Une phrase
entre crochets, « [Il suffit] bien sûr … », est en partie illisible.}.

\end{document}
